% SGA 3，第 VII_B 篇，英文重建稿——第 1 轮。
% 范围：权威底本分篇第 78--79 页，第 4.3.2 节全文，止于第 4.4 节之前。
% 合并阅读本页码： 591--592.
% 再版印刷页码： 580--581.
% 本范围内原版页边定位： 548.
% 本范围内编者注： 138.
% 权威底本： Polo--Gille Exp7B-13oct24.pdf,
% SHA-256 D303690E0195654BBE7C3DAA5D89EA23E24CD709A3575217234AC19AE8836694.
% 曾参考 Jacob Reinhold 的 CC BY 4.0 Markdown 作为比对谱系。
% 原英文 Loop 1 阶段说明：自然性方图采用直接取自 PDF 的精确裁片。当前中文工程已采用原生图。

% 2026-10-22：对照英文底本逐段复核；保留原有译文，核对公式、引用键、脚注与原生图路径。
\subsubsection*{4.3.2}
\phantomsection
\label{sga3:VIIB:4.3.2}

“在特征 \(p>0\) 时不存在 Campbell--Hausdorff 公式。”说明如下。\SGARefLink{sga3:VIIA:4.3.1}{4.3.1} 的函子性同构依赖于基
\(({}^C x_i)_{i\in I}\) 的选择。与特征 \(0\) 时 \SGARefLink{sga3:VIIB:3.2:prop}{3.2} 中的情形不同，一般不存在“同时关于 \(C\) 和 \(\mathbf L\) 函子性”的双射
\[
  \mathbf L(C)\otimes_C\sqrt[p]{0_C}
  \longrightarrow
  \mathscr G_p(\mathbf L)(C)
\]
。例如，下面的论证表明，当 \(k\) 是一个包含所有
\((p^2-1)\) 次单位根的域时，不存在这样的同构。

选取 \(\mathbf L\)，使得对每个
\(C\in\operatorname{Alf}/k\)，代数 \(\mathbf L(C)\) 为由一个元素 \(x\) 生成、满足关系
\[
  x^{(p^2)}=0.
\]
的阿贝尔 \(C\) 上 \(p\)-Lie代数。因而
\[
  \mathbf L(C)
  =
  Cx\oplus Cx^{(p)},
  \qquad
  U_p\bigl(\mathbf L(C)\bigr)
  \simeq
  C[x]/(x^{p^2}).
\]
\phantomsection
\label{sga3:VIIB:4.3.2:source-notation-note}
\noindent\textit{底本说明。}\enspace
法文原文在上式中印作 \(k[x]/(x^{p^2})\)，但这里的包络代数是在 \(C\) 上构造的。稍后原文又将由群状元素组成的 \(\mathscr G_p(\mathbf L)(C)\) 与
\(\operatorname{Prim}U_p(\mathbf L(C))\) 等同。本底本使用系数环 \(C\)，并保留 \SGARefLink{sga3:VIIA:4.2.1}{4.2.1} 中定义的群状元素集合；本原元素则在 4.2.2 中描述它的 Lie代数。

我们将证明，与 \(\mathbf L\) 的自同态相容、并将
\(\mathbf L(C)\otimes_C\sqrt[p]{0_C}\) 映入
\(\mathscr G_p(\mathbf L)(C)\) 的唯一函子性态射
\[
  \chi_C:
  \mathbf L(C)\otimes_C\sqrt[p]{0_C}
  \longrightarrow
  U_p\bigl(\mathbf L(C)\bigr)
\]
是取值为 \(1\) 的常值态射。

设
\[
  \psi_C:
  \mathbf L(C)\otimes_C\sqrt[p]{0_C}
  \xrightarrow{\ \sim\ }
  \mathscr G_p(\mathbf L)(C)
\]
为 \SGARefLink{sga3:VIIA:4.3.1}{4.3.1} 的双射，即
\[
  x\otimes a+x^{(p)}\otimes b
  \longmapsto
  \sum_{0\leq i,j<p}
  a^i b^j x^{\,i+pj}.
\]
将 \(\chi_C\) 与 \(\psi_C^{-1}\) 复合，得到关于 \(C\) 函子性的态射
\[
\begin{aligned}
  \theta_C:
  \mathbf L(C)\otimes_C\sqrt[p]{0_C}
  &\longrightarrow
  \mathbf L(C)\otimes_C\sqrt[p]{0_C},
  \\
  x\otimes a+x^{(p)}\otimes b
  &\longmapsto
  x\otimes P(a,b)+x^{(p)}\otimes Q(a,b).
\end{aligned}
\]
关于 \(C\) 的函子性蕴含 \(P(a,b)\) 与 \(Q(a,b)\) 是两个多项式
\[
  P,Q\in k[X,Y]
\]
在 \((a,b)\) 处的值，并可假定这两个多项式关于 \(X\) 和 \(Y\) 中每一个的次数都小于 \(p\)。\footnote{\label{sga3:VIIB:nde:138}编者扩充了原文接下来的论证。}

设 \(\alpha\in k\)，令 \(\ell_\alpha\) 为 \(p\)-Lie代数 \(\mathbf L\) 的如下自同态：
\[
  x\longmapsto\alpha x,
  \qquad
  x^{(p)}\longmapsto\alpha^p x^{(p)}.
\]
则 \(U_p(\ell_\alpha)\) 将 \(x\) 送至 \(\alpha x\)，因而对 \(i<p^2\)，将 \(x^i\) 送至 \(\alpha^i x^i\)。下列方图交换：
\begin{center}
  \phantomsection
  \label{sga3:VIIB:4.3.2:naturality-endomorphism-square}
  % SGA3-VII-DIAG-123
  % SGA3-VII-DIAG-123
% Authority: Exp7B-13oct24.pdf, local p. 78.
% Crop: figures/Exp7B_localp078_naturality_endomorphism_square.png
% Locator: tex/components/71_expose_VIIB_section432_no_functorial_campbell_hausdorff_en.tex, Loop 1 line 121.
\begin{tikzcd}[row sep=3.2em,column sep=4.8em]
\mathbf L(C)\otimes_C\sqrt[p]{0_C}
  \arrow[r,"\psi_C"]
  \arrow[d,"{\ell_\alpha(C)\otimes_C\operatorname{id}}"]
  & U_p(\mathbf L(C))
      \arrow[d,"{U_p(\ell_\alpha)(C)}"] \\
\mathbf L(C)\otimes_C\sqrt[p]{0_C}
  \arrow[r,"\psi_C"]
  & U_p(\mathbf L(C))
\end{tikzcd}

\end{center}
因此，假设
\[
  \chi_C\circ(\ell_\alpha(C)\otimes\operatorname{id})
  =
  U_p(\ell_\alpha)(C)\circ\chi_C
\]
给出
\[
  P(\alpha a,\alpha^p b)=\alpha P(a,b),
  \qquad
  Q(\alpha a,\alpha^p b)=\alpha^p Q(a,b).
\]

写成
\[
  P=\sum_{i,j<p}\lambda_{ij}X^iY^j,
  \qquad
  Q=\sum_{i,j<p}\mu_{ij}X^iY^j,
\]
并取
\[
  C=k[X,Y]/(X^p,Y^p).
\]
若 \(\lambda_{ij}\ne0\)，则
\[
  \alpha^{\,i-1+pj}=1;
\]
若 \(\mu_{ij}\ne0\)，则
\[
  \alpha^{\,i+p(j-1)}=1.
\]
取 \(\alpha\) 为一个 \(p^2-1\) 阶本原单位根，就得到
\[
  P=\lambda X,
  \qquad
  Q=\mu Y
\]
其中 \(\lambda,\mu\in k\)。因此，
\[
  \chi_C(x\otimes a+x^{(p)}\otimes b)
  =
  \sum_{0\leq i,j<p}
  (\lambda a)^i(\mu b)^j x^{\,i+pj}.
\]

最后，考虑 \(\mathbf L\) 的如下自同态 \(f\)：
\[
  x\longmapsto x+x^{(p)}.
\]
取 \(C=k[a]/(a^3)\)，比较
\[
  (\chi_C\circ f(C))(x\otimes a)
  \qquad\text{和}\qquad
  \bigl(U_p(f)(C)\circ\chi_C\bigr)(x\otimes a).
\]
中 \(x^p\) 与 \(x^{p+1}\) 的系数。这给出
\[
  \lambda=\mu,
  \qquad
  \lambda^2a^2=0,
\]
从而 \(\lambda=\mu=0\)，得证。
